\section{Discussion}
\label{sec:discussion}

The design choices in Frankenstein Transformer reflect several engineering and research tensions in modern deep learning tooling.

The campaign reported in Section~\ref{sec:empirical-evaluation} turns several of those tensions into measurements rather than opinions. The strict schema paid off operationally: 868 configurations trained, deployed and were evaluated through the same contract, and the failures that remained were concentrated in four mixers and in checkpoint serialisation, not in configuration drift. The prefixed optimizer system proved expressive but, at 100 epochs, undecisive---a dozen optimizers converged to the same loss---which is itself an argument for the schema: only a controlled sweep can show that a design dimension does not matter at a given scale.

\subsection{Schema-Driven Design Trade-offs}

The schema-first approach provides significant reproducibility benefits by enforcing explicit contracts and failing fast on invalid configurations. However, this approach also introduces rigidity: adding new architectures or optimizers requires schema extensions rather than loose command-line arguments. The prefixed hyperparameter system enables fine-grained control but increases configuration complexity for users accustomed to simpler interfaces.

The decision to enforce \texttt{additionalProperties: false} at all schema levels eliminates silent parameter swallowing that has plagued earlier configuration systems, but this strictness requires careful schema maintenance when extending system capabilities. Each new attention mechanism or optimizer variant must be properly integrated into the validation framework, including schema field definitions with appropriate types and constraints, prefixed hyperparameter mapping for optimizer-specific groups, default values aligned with research best practices, and documentation strings for web interface rendering.

\subsection{Architectural Coverage and Gaps}

The forty-nine implemented mixers span the major research directions in sequence modeling, but certain gaps remain. The system lacks recent hybrid architectures such as Griffin \citep{soares_griffin_2024} and Jamba \citep{gu_jamba_2024}, which combine gating with state-space models. MoE routing is implemented for FFN layers but not for attention computation, where recent work has shown benefits \citep{lewis_moe_2021}.

The sparse attention coverage is broad. The two inference-time methods (FASA, SpargeAttn) keep their training-free character in the sense that their sparsification rule is never learned; instead, the implementation trains the layer as standard dense attention over the same projections and activates the sparse selection/pruning path only at evaluation, so both mixers can be trained from scratch and their sparse behaviour is exactly the one described in the original papers.

The gated and fast-weight coverage spans the major categories, including the six Falcon variants \citep{zhang_falcon_2026}. Residual-stream macro-design is covered by Attention Residuals \citep{sun_attnres_2026} and by the manifold-constrained $n$-stream residual (mHC) \citep{xie_mhc_2025}, whose Sinkhorn--Knopp projection onto the Birkhoff polytope restores the identity-mapping property lost by unconstrained hyper-connections; mHC is validated only at the level of the published 3B--27B MoE experiments and would benefit from independent reproductions at smaller scales (Appendix~\ref{sec:annex-mhc}).

\subsection{Optimizer Landscape Fragmentation}

The support for thirty-eight optimizer families across eight categories demonstrates comprehensiveness but also highlights the fragmented state of optimization research. Users face significant decision complexity when choosing among variance-reduction methods (Adan, MARS), memory-efficient variants (GaLore, Adafactor, APOLLO), and curvature-aware approaches (Sophia, Shampoo). The prefixed hyperparameter system, while powerful, requires understanding of which parameters are relevant for each optimizer class.

The implementation quality varies across optimizers: classical methods (AdamW, SGD with momentum) are highly optimized in PyTorch, while newer methods (Muon, Turbo-Muon, AF-Muon, SOAP) may require custom implementations that affect numerical stability and performance characteristics.

\subsection{Deployment and Production Considerations}

The quantization pipeline demonstrates practical deployment concerns but makes specific engineering trade-offs. Ternary weight packing reduces storage to approximately $1.58$ bits per parameter, but this aggressive compression may degrade performance, especially for smaller models where quantization error is more significant. The safetensors exporter instead offers block-affine \texttt{Q8\_0} and \texttt{Q4\_0} formats (8.5 and 4.5 bits per weight) and applies all three only to two-dimensional linear weights whose input dimension is divisible by 32; embeddings, normalization parameters and biases remain in FP32, so the effective compression ratio depends on the linear fraction of the model. The packed-linear runtime keeps the artifact compact at the cost of a dequantization step in every forward pass.

The SBERT workflows provide practical utility for semantic similarity and retrieval tasks, but implementation assumes standard pooling strategies (CLS token, mean pooling). Recent advances such as Matryoshka embeddings \citep{muennighoff_matryoshka_2021} and contrastive learning refinements \citep{chen_supervised_2020} are not yet incorporated.

\subsection{Integration and Extensibility Challenges}

Extensibility is handled by a declarative component registry rather than ad-hoc dispatch. Every pluggable component---sequence mixer, optimizer, normalization, activation, residual strategy, positional encoding and token embedding---is described by a frozen \texttt{ComponentSpec} registered through \path{src/registry/}, and the schema enumerations, per-component parameter fragments, optimizer prefix reference and conditional prefix rules are assembled from that registry before validation (Section~\ref{sec:system-design}). A third-party plugin registers the same kind of spec through the public authoring API, either from the configuration's \texttt{plugins:} list or through the \texttt{frankenstein.plugins} entry-point group, so it validates strictly, appears in the schema-generated builder and is dispatched by the same core code as the built-ins. The maintenance pressure is therefore the size of the component catalog itself---forty-nine mixers and thirty-eight optimizers---rather than the dispatch logic, which contains no per-optimizer branches and no per-mixer selection ladder.

The web configuration interface provides significant usability improvements but introduces deployment complexity: running Streamlit alongside training jobs requires additional resources and infrastructure considerations that may not be appropriate for all environments, particularly HPC clusters without web access.

